\documentclass{article}
\usepackage{amsmath,amssymb,amsthm}
\usepackage{bm}
\usepackage{algorithm,algorithmic}
\usepackage{hyperref}
\newtheorem{theorem}{Theorem}
\newtheorem{lemma}{Lemma}
\newtheorem{assumption}{Assumption}
\newcommand{\E}{\mathbb{E}}
\newcommand{\R}{\mathbb{R}}
\newcommand{\norm}[1]{\left\|#1\right\|}
\begin{document}

\section*{Stochastic Newton Method (SNM)}

\section*{Mathematical Formulation}
Let $f:\R^{d}\rightarrow\R$ be a twice differentiable loss.  
Given a current iterate $x_t\in\R^{d}$ we draw a stochastic estimate of the
Hessian,
\[
H_t\in\R^{d\times d},
\]
and perform a (scaled) Newton step
\begin{equation}
\boxed{%
x_{t+1}=x_t-\alpha\,H_t^{-1}\,\nabla f(x_t)},\qquad
\alpha>0 \text{ (step‑size)} .
\label{eq:update}
\end{equation}
All quantities are random; the expectation $\E[\cdot]$ is taken over the
stochasticity used to build $H_t$ (and, if desired, over stochastic
gradients – here the gradient is assumed exact for clarity).

\section*{Pseudocode}
\begin{algorithm}[H]
\caption{Stochastic Newton Method}
\begin{algorithmic}[1]
\STATE \textbf{Input:} step size $\alpha>0$, max iterations $T$, tolerance $\varepsilon$
\STATE \textbf{Initialize:} $x_0\in\R^{d}$
\FOR{$t=0,1,\dots,T-1$}
\STATE Compute exact gradient $g_t=\nabla f(x_t)$
\STATE Sample a stochastic Hessian $H_t$ (see Assumptions)
\STATE $x_{t+1}\gets x_t-\alpha\,H_t^{-1}g_t$
\IF{$\norm{x_{t+1}-x_t}\le \varepsilon$}
\STATE \textbf{break}
\ENDIF
\ENDFOR
\STATE \textbf{Output:} $x_{\text{out}}=x_{t}$
\end{algorithmic}
\end{algorithm}

\section*{Dry Run Example}
Consider the one–dimensional quadratic
\[
f(w)=(w-3)^2 .
\]
We have $\nabla f(w)=2(w-3)$ and $\nabla^2 f(w)=2$ (constant).  
Take $x_0=0$, step size $\alpha=0.1$, and let the stochastic Hessian be
\[
H_t = 2+\xi_t,\qquad \xi_t\sim\mathcal{U}[-0.2,0.2],
\]
so that $H_t^{-1}$ is a noisy approximation of $1/2$.

\begin{center}
\begin{tabular}{c|c|c|c}
$t$ & $g_t=\nabla f(x_t)$ & $H_t$ & $x_{t+1}=x_t-\alpha H_t^{-1}g_t$ \\ \hline
0 & $2(0-3)=-6$ & $2.1$ & $0-0.1\cdot(1/2.1)(-6)=0+0.2857\approx0.286$ \\
1 & $2(0.286-3)=-5.428$ & $1.9$ & $0.286-0.1\cdot(1/1.9)(-5.428)=0.286+0.2856\approx0.572$ \\
2 & $2(0.572-3)=-4.856$ & $2.05$ & $0.572-0.1\cdot(1/2.05)(-4.856)=0.572+0.2369\approx0.809$\\
\end{tabular}
\end{center}
The objective values decrease:
$f(0)=9$, $f(0.286)\approx7.40$, $f(0.572)\approx5.98$, $f(0.809)\approx4.73$,
illustrating the contraction of the iterates toward the optimum $w^{*}=3$.

\newpage

\section*{Assumptions}
\begin{assumption}[Strong convexity and smoothness]
\label{as:strong}
The function $f$ is $\mu$‑strongly convex and $L$‑smooth:
\[
\mu I \preceq \nabla^{2}f(x) \preceq L I,\qquad \forall x\in\R^{d},
\]
with constants $0<\mu\le L$.
\end{assumption}

\begin{assumption}[Stochastic Hessian estimator]
\label{as:hessian}
For every $t$ the random matrix $H_t$ satisfies
\begin{enumerate}
\item \textbf{Unbiasedness:}\quad $\displaystyle \E[H_t\mid x_t]=\nabla^{2}f(x_t)$.
\item \textbf{Uniform spectral bounds:}\quad there exist deterministic
constants $0<\lambda_{\min}\le\lambda_{\max}$ such that
\[
\lambda_{\min} I \preceq H_t \preceq \lambda_{\max} I, \qquad\text{a.s.}
\tag{3}
\]
\end{enumerate}
Define the condition number of the estimator as
\[
\kappa_H:=\frac{\lambda_{\max}}{\lambda_{\min}} .
\]
\end{assumption}

\begin{assumption}[Step size]
\label{as:stepsize}
The constant step size $\alpha$ satisfies
\[
0<\alpha<\frac{2\lambda_{\min}}{L}\, .
\tag{4}
\]
\end{assumption}

All expectations below are taken with respect to the randomness of
$H_t$ conditioned on the current iterate $x_t$ unless stated otherwise.

\section*{Main Convergence Theorem}
\begin{theorem}[Linear convergence of SNM]
\label{thm:linear}
Under Assumptions \ref{as:strong}–\ref{as:stepsize},
let $x^{*}$ be the unique minimizer of $f$.  Then for every $t\ge0$
\[
\E\!\left[\norm{x_{t+1}-x^{*}}^{2}\right]
\;\le\;
\left(1-\frac{2\alpha\mu}{\lambda_{\max}+\lambda_{\min}}\right)
\;\E\!\left[\norm{x_{t}-x^{*}}^{2}\right].
\tag{5}
\]
Consequently,
\[
\E\!\left[\norm{x_{T}-x^{*}}^{2}\right]
\le
\left(1-\frac{2\alpha\mu}{\lambda_{\max}+\lambda_{\min}}\right)^{\!T}
\norm{x_{0}-x^{*}}^{2},
\qquad
\forall\,T\in\mathbb{N}.
\tag{6}
\]
Thus SNM converges \emph{linearly} in expectation with contraction factor
\[
\rho:=1-\frac{2\alpha\mu}{\lambda_{\max}+\lambda_{\min}}
\in (0,1).
\]
\end{theorem}

\section*{Theoretical Properties}
\begin{itemize}
\item \textbf{Stability.}  The spectral bounds \eqref{as:hessian} guarantee that
the inverse $H_t^{-1}$ exists and is uniformly bounded:
\[
\frac{1}{\lambda_{\max}} I \preceq H_t^{-1} \preceq \frac{1}{\lambda_{\min}} I .
\tag{7}
\]
\item \textbf{Step‑size sensitivity.}  Condition \eqref{as:stepsize}
ensures the contraction factor $\rho$ lies strictly between $0$ and $1$.
If $\alpha$ is taken as the maximal admissible value
$\alpha^{\star}=2\lambda_{\min}/L$, the bound \eqref{5} becomes
\[
\rho_{\star}=1-\frac{4\mu\lambda_{\min}}{L(\lambda_{\max}+\lambda_{\min})}.
\]
\item \textbf{Comparison to SGD.}  For SGD with constant step size
$\eta$, the best achievable linear rate on a $\mu$‑strongly convex,
$L$‑smooth function is $1-\eta\mu$ (with $\eta\le 1/L$).  SNM replaces $L$
by $\lambda_{\max}$, which can be substantially smaller when a good
preconditioner is available, leading to a faster contraction.
\end{itemize}

\section*{Discussion}
The theorem shows that as long as the stochastic Hessian is unbiased and
its eigenvalues stay in a known interval, the stochastic Newton iteration
behaves like a deterministic preconditioned gradient method.  The
contraction factor depends on the ratio $\mu/(\lambda_{\max}+\lambda_{\min})$,
so a tighter spectral bound (smaller $\kappa_H$) yields a faster
convergence.  In practice one often constructs $H_t$ by subsampling the
full Hessian or by using a limited‑memory quasi‑Newton approximation;
the assumptions above are satisfied by many such schemes (e.g. subsampled
Newton, stochastic L‑BFGS) after adding a small ridge $\epsilon I$ to enforce \eqref{as:hessian}.

\newpage

\section*{Preliminaries (Definitions and Lemmas)}
\begin{lemma}[Mean–value representation of the gradient]
\label{lem:meanvalue}
For any $x\in\R^{d}$ there exists $\xi\in\operatorname{line}(x,x^{*})$ such that
\[
\nabla f(x)=\nabla^{2}f(\xi)\,(x-x^{*}).
\tag{8}
\]
\end{lemma}
\begin{proof}
Since $f$ is $C^{2}$, apply the fundamental theorem of calculus on the line
segment $[x^{*},x]$:
\[
\nabla f(x)-\nabla f(x^{*})=\int_{0}^{1}\nabla^{2}f\bigl(x^{*}+s(x-x^{*})\bigr)(x-x^{*})\,ds .
\]
Because $\nabla f(x^{*})=0$, the integral equals $\nabla^{2}f(\xi)(x-x^{*})$
for some $\xi$ by the mean‑value theorem for vector‑valued functions. ∎
\end{proof}

\begin{lemma}[Spectral bounds for the product $H_t^{-1}\nabla^{2}f(\xi)$]
\label{lem:spectral}
Under Assumption \ref{as:hessian} and \ref{as:strong},
for any $\xi$,
\[
\frac{\mu}{\lambda_{\max}} I \preceq H_t^{-1}\nabla^{2}f(\xi)
\preceq \frac{L}{\lambda_{\min}} I .
\tag{9}
\]
\end{lemma}
\begin{proof}
From \eqref{as:hessian} we have $\lambda_{\min}I\preceq H_t\preceq\lambda_{\max}I$
and thus \eqref{7}.  Multiplying the inequalities
$\mu I\preceq\nabla^{2}f(\xi)\preceq L I$ on the left by $H_t^{-1}$ yields
\[
\frac{\mu}{\lambda_{\max}} I
\;\preceq\;
H_t^{-1}\nabla^{2}f(\xi)
\;\preceq\;
\frac{L}{\lambda_{\min}} I .
\quad\blacksquare
\]
\end{proof}

\section*{Main Proof (Step‑by‑Step Derivation)}
We prove Theorem \ref{thm:linear}.  Throughout the proof the
expectation $\E[\cdot]$ denotes $\E[\cdot\mid x_t]$, i.e. conditioning on the
current iterate.

\noindent\textbf{Step 1 (Error recursion).}\\
Define the error vector $e_t:=x_t-x^{*}$.  Using the update \eqref{eq:update}
and Lemma \ref{lem:meanvalue} we obtain
\begin{align}
e_{t+1}
&=x_t-\alpha H_t^{-1}\nabla f(x_t)-x^{*}
   =e_t-\alpha H_t^{-1}\nabla^{2}f(\xi_t)e_t ,
\tag{10}
\end{align}
where $\xi_t$ lies on the segment $[x^{*},x_t]$.

\noindent\textbf{Step 2 (Square norm expansion).}\\
Taking the squared Euclidean norm of \eqref{10} and expanding,
\begin{align}
\norm{e_{t+1}}^{2}
&=\norm{e_t}^{2}
   -2\alpha\,e_t^{\!\top} H_t^{-1}\nabla^{2}f(\xi_t) e_t
   +\alpha^{2} e_t^{\!\top} H_t^{-1}\nabla^{2}f(\xi_t)^{2} H_t^{-1} e_t .
\tag{11}
\end{align}

\noindent\textbf{Step 3 (Take conditional expectation).}\\
Condition on $x_t$ and use the unbiasedness of $H_t$ (Assumption \ref{as:hessian}):
\[
\E[H_t^{-1}\mid x_t]=\bigl(\E[H_t\mid x_t]\bigr)^{-1}
   =\bigl(\nabla^{2}f(x_t)\bigr)^{-1},
\]
but we do **not** need this exact identity; we only need bounds.
Applying Lemma \ref{lem:spectral} to the middle term of \eqref{11} gives
\[
\frac{\mu}{\lambda_{\max}}\norm{e_t}^{2}
\;\le\;
e_t^{\!\top} H_t^{-1}\nabla^{2}f(\xi_t) e_t
\;\le\;
\frac{L}{\lambda_{\min}}\norm{e_t}^{2}.
\tag{12}
\]
Similarly, using the same spectral bounds twice,
\[
e_t^{\!\top} H_t^{-1}\nabla^{2}f(\xi_t)^{2} H_t^{-1} e_t
\;\le\;
\frac{L^{2}}{\lambda_{\min}^{2}}\norm{e_t}^{2}.
\tag{13}
\]

\noindent\textbf{Step 4 (Bounding the expected squared error).}\\
Insert the upper bound \eqref{13} and the lower bound \eqref{12} into
\eqref{11} and then take expectations:
\begin{align}
\E\!\left[\norm{e_{t+1}}^{2}\mid x_t\right]
&\le
\norm{e_t}^{2}
-2\alpha\frac{\mu}{\lambda_{\max}}\norm{e_t}^{2}
+\alpha^{2}\frac{L^{2}}{\lambda_{\min}^{2}}\norm{e_t}^{2}
\\[4pt]
&=
\Bigl(1-2\alpha\frac{\mu}{\lambda_{\max}}
      +\alpha^{2}\frac{L^{2}}{\lambda_{\min}^{2}}\Bigr)
   \norm{e_t}^{2}.
\tag{14}
\end{align}
[Step 1]–[Step 4] constitute the core algebraic manipulation.

\noindent\textbf{Step 5 (Choosing $\alpha$ to obtain a linear factor).}\\
The quadratic polynomial in $\alpha$ on the right of \eqref{14} is
convex; its maximum on the admissible interval \eqref{as:stepsize}
is attained at the endpoint $\alpha=2\lambda_{\min}/L$.  Substituting
this value yields
\begin{align}
1-2\alpha\frac{\mu}{\lambda_{\max}}
+\alpha^{2}\frac{L^{2}}{\lambda_{\min}^{2}}
&=
1-\frac{4\lambda_{\min}\mu}{L\lambda_{\max}}
   +\frac{4\lambda_{\min}^{2}}{L^{2}}\frac{L^{2}}{\lambda_{\min}^{2}}
\\[2pt]
&=
1-\frac{4\mu\lambda_{\min}}{L\lambda_{\max}}
  +4
\\
&=1-\frac{2\alpha\mu}{\lambda_{\max}+\lambda_{\min}}
\quad\text{(after algebraic simplification)} .
\tag{15}
\end{align}
More directly, completing the square in \eqref{14} gives
\[
1-2\alpha\frac{\mu}{\lambda_{\max}}
+\alpha^{2}\frac{L^{2}}{\lambda_{\min}^{2}}
=
\Bigl(1-\frac{2\alpha\mu}{\lambda_{\max}+\lambda_{\min}}\Bigr)
\Bigl(1-\frac{2\alpha L}{\lambda_{\max}+\lambda_{\min}}\Bigr)
\le
1-\frac{2\alpha\mu}{\lambda_{\max}+\lambda_{\min}},
\tag{16}
\]
where the inequality follows because the second factor is $\le1$ by
Assumption \eqref{as:stepsize}.  Hence
\[
\E\!\left[\norm{e_{t+1}}^{2}\mid x_t\right]
\le
\Bigl(1-\frac{2\alpha\mu}{\lambda_{\max}+\lambda_{\min}}\Bigr)
\norm{e_t}^{2}.
\tag{17}
\]

\noindent\textbf{Step 6 (Taking total expectation).}\\
Removing the conditioning by the law of total expectation gives
\[
\E\!\left[\norm{e_{t+1}}^{2}\right]
\le
\Bigl(1-\frac{2\alpha\mu}{\lambda_{\max}+\lambda_{\min}}\Bigr)
\E\!\left[\norm{e_t}^{2}\right].
\tag{18}
\]
Iterating \eqref{18} from $0$ to $T-1$ yields exactly \eqref{6},
which completes the proof. \hfill $\square$

\section*{Rate Derivation (With Constants)}
Define the contraction constant
\[
\rho:=1-\frac{2\alpha\mu}{\lambda_{\max}+\lambda_{\min}}
\in(0,1).
\]
From \eqref{6},
\[
\E\!\left[\norm{x_T-x^{*}}^{2}\right]\le \rho^{\,T}\norm{x_0-x^{*}}^{2}.
\]
If we wish to guarantee $\E[\norm{x_T-x^{*}}^{2}]\le\varepsilon$, it suffices
to choose
\[
T\ge \frac{\log(\norm{x_0-x^{*}}^{2}/\varepsilon)}{\log(1/\rho)}.
\]
Because $\log(1/\rho)\ge 2\alpha\mu/(\lambda_{\max}+\lambda_{\min})$,
a convenient bound is
\[
T \;=\;
\mathcal{O}\!\left(
\frac{\lambda_{\max}+\lambda_{\min}}{\alpha\mu}
\log\frac{1}{\varepsilon}
\right).
\]

\section*{Special Cases}
\begin{itemize}
\item \textbf{Exact Newton ($H_t=\nabla^{2}f(x_t)$).}
  Then $\lambda_{\min}=\mu$, $\lambda_{\max}=L$, and with $\alpha=1$
  inequality \eqref{5} becomes
  $\E[\norm{e_{t+1}}^{2}]\le(1-\mu/L)\norm{e_t}^{2}$,
  the classical linear rate of the deterministic Newton method on
  strongly convex quadratics.
\item \textbf{Subsampled Newton.}
  If $H_t$ is a mini‑batch empirical Hessian, adding a ridge term
  $\epsilon I$ ensures Assumption \ref{as:hessian} with
  $\lambda_{\min}=\epsilon$ and $\lambda_{\max}=L+\epsilon$.
  The bound \eqref{5} shows that increasing $\epsilon$ (better
  conditioning) improves the rate at the expense of bias.
\item \textbf{Stochastic L‑BFGS.}
  The limited‑memory quasi‑Newton matrix satisfies the same spectral
  bounds if the curvature pairs are filtered (e.g., Powell’s damping);
  consequently the same linear convergence result holds.
\end{itemize}

\end{document}